%! Adding \& Subtracting Rational Expressions — Unit 4, Lesson 4.3 — Warm-Up (Answer Key)
\documentclass[10pt]{article}
\usepackage{algebra2-article}
\usepackage{algebra2-key}   % pulls in -boxes; do NOT also load -boxes
\newcommand{\TallMath}[1]{$\displaystyle #1\rule[-1.4em]{0pt}{3.2em}$}

\begin{document}

\pageheader{Unit 4, Lesson 4.3}{Warm-Up --- Answer Key}
\namedateperiod

\vspace{0.10in}

\begin{notesbox}{Getting Ready: Skills We'll Use Today}
\small Three skills, one per leg of today's lesson. Work quickly.

\vspace{4pt}
\textbf{1. Number fractions --- the common denominator.}
\begin{enumerate}[label=(\alph*), leftmargin=2.1em, itemsep=3pt]
  \item Write each denominator as a product of primes, then build the \emph{smallest} denominator
        that both can become. \\[2pt]
        $12=$ \ans{$2^2\cdot3$}, \quad $18=$ \ans{$2\cdot3^2$}, \quad so the least common denominator is
        \ans{36}. \\[3pt]
        $\dfrac{5}{12}+\dfrac{7}{18}=$ \ans{$\frac{15}{36}$} $+$ \ans{$\frac{14}{36}$} $=$
        \ans{$\frac{29}{36}$}
  \item You could also have used $12\cdot18=216$ as the denominator. Rewrite both fractions over
        $216$ and add: \\[3pt]
        \ans{$\frac{90}{216}$} $+$ \ans{$\frac{84}{216}$} $=$ \ans{$\frac{174}{216}$} \\[3pt]
        Same value as part (a)? \ans{yes} \quad What extra work did it cost you?
        \ans{reducing $\frac{174}{216}$ back to $\frac{29}{36}$}
  \item To \emph{divide} by a fraction, multiply by its \ans{reciprocal} (Lesson 4.2):
        $\dfrac{2}{3}\div\dfrac{5}{6}=\dfrac{2}{3}\cdot$ \ans{$\frac{6}{5}$} $=$ \ans{$\frac{4}{5}$}
\end{enumerate}

\vspace{4pt}
\textbf{2. Signs and factors.}
\begin{enumerate}[label=(\alph*), leftmargin=2.1em, itemsep=3pt]
  \item $9-(4-6)=$ \ans{11} \qquad but \qquad $9-4-6=$ \ans{$-1$}. \\[2pt]
        They are not equal, because a minus sign in front of a parenthesis changes
        \ans{the sign of every term} inside it.
  \item Factor completely (the Unit 3 toolkit): \; $x^2-9=$ \ans{$(x-3)(x+3)$}; \quad
        $x^2-4=$ \ans{$(x-2)(x+2)$}; \quad $x^2+5x+6=$ \ans{$(x+2)(x+3)$}
\end{enumerate}

\vspace{4pt}
\textbf{3. Simplify and restrict} (Lesson 4.1). For $\dfrac{3x-9}{x^2-9}$:
\begin{enumerate}[label=(\alph*), leftmargin=2.1em, itemsep=3pt]
  \item Factored form: \ans{$\frac{3(x-3)}{(x-3)(x+3)}$} \quad Restrictions: $x\neq$ \ans{$3,\ -3$}
  \item Simplest form: \ans{$\frac{3}{x+3}$}
  \item Which excluded value is a \emph{hole}? \ans{$x=3$} \; Which is a \emph{wall} (vertical
        asymptote)? \ans{$x=-3$}
\end{enumerate}

\end{notesbox}

\begin{teachernote}
Each item is one leg of today's lesson. \textbf{Item 1(a)} is the whole procedure of the day done on
numbers, where students already trust it: \emph{factor the denominators, build the LCD out of those
factors, rewrite, then combine}. \textbf{1(b)} answers the question every class asks --- ``why not just
multiply the denominators?'' --- honestly: you may, and you will get the same value, but you buy
yourself a reduction at the end. That is exactly the argument for the \emph{least} common denominator
with polynomials, where the reduction at the end is a re-factoring job. \textbf{1(c)} keeps
keep--change--flip warm; you need it in Section 4, because a complex fraction \emph{is} a division.

\vspace{3pt}
\textbf{Item 2(a)} is the seed of today's signature error, planted with numbers before any letters
appear. Keep it on the board. When a student writes $\frac{4x-1-2x-7}{x+3}$ in the Hook, do not quote a
rule --- point at $9-(4-6)$.

\vspace{3pt}
\textbf{Item 3} is the payoff to watch. Its answer, $\frac{3}{x+3}$ with $x\neq3,-3$, is the
\emph{same} answer today's anchor problem produces from a completely different starting point
($\frac{3}{x-3}-\frac{18}{x^2-9}$). Say nothing about that now; when it lands in Section 3, ask who
recognizes it.
\end{teachernote}

\end{document}
